\documentclass[10pt]{article}
\usepackage[T1]{fontenc}
\usepackage[utf8]{inputenc}
\usepackage[french]{babel}
\usepackage{amsmath,amssymb,mathrsfs,mathpazo}
\usepackage{geometry}
\geometry{a4paper,margin=2cm}
\usepackage{graphicx,array,tabularx,tabularray,colortbl,multirow,multicol,booktabs,enumitem,pifont,ulem,fancybox,diagbox,cancel}
\usepackage{pstricks,pstricks-add,pst-plot,pst-tree,pst-node,pst-text,pst-func,pst-math,pst-bspline,pst-3dplot,pst-eucl,pst-solides3d}
\let\pscalculate\undefined
\usepackage{pst-bezier}
\usepackage{tikz,pgfplots,tkz-tab}
\usetikzlibrary{arrows,arrows.meta,patterns,calc,positioning,trees,decorations.markings,shapes.geometric,shapes.misc,shapes.arrows,intersections,decorations.pathmorphing,decorations.pathreplacing,backgrounds,shadows,babel}
\pgfplotsset{compat=1.18}
\usepgfplotslibrary{fillbetween}
\usepackage[np]{numprint}
\usepackage{listings,makecell,pas-tableur,twoopt}
\let\barre\undefined
\usepackage[pstricks]{tablvar}
\definecolor{darkgreen}{RGB}{0,100,0}
\providecommand{\red}{\color{red}}
\providecommand{\blue}{\color{blue}}
\providecommand{\green}{\color{green}}
\providecommand{\black}{\color{black}}
\providecommand{\white}{\color{white}}
\providecommand{\gray}{\color{gray}}
\providecommand{\cyan}{\color{cyan}}
\providecommand{\magenta}{\color{magenta}}
\providecommand{\yellow}{\color{yellow}}
\providecommand{\violet}{\color{violet}}
\providecommand{\orange}{\color{orange}}
\providecommand{\euro}{\texteuro}
\providecommand{\EUR}{\texteuro}
\providecommand{\ieme}{\textsuperscript{e}}
\providecommand{\degres}{\ensuremath{^\circ}}
\providecommand{\degree}{\ensuremath{^\circ}}
\providecommand{\vv}[1]{\overrightarrow{#1}}
\providecommand{\Coord}[1]{#1}
\providecommand{\decosix}{}
\providecommand{\starredbullet}{\ensuremath{\bullet}}
\pagestyle{empty}
\setlength{\parindent}{0pt}
\usepackage{adjustbox}\DeclareUnicodeCharacter{2236}{:}\DeclareUnicodeCharacter{2190}{\ensuremath{\leftarrow}}
\providecommand{\R}{\mathbb{R}}
\providecommand{\N}{\mathbb{N}}
\providecommand{\Z}{\mathbb{Z}}
\providecommand{\Q}{\mathbb{Q}}
\providecommand{\D}{\mathbb{D}}
\providecommand{\C}{\mathbb{C}}
\providecommand{\e}{{\mathrm{e}}}
\providecommand{\vect}[1]{\overrightarrow{#1}}
\providecommand{\vectt}[1]{\overrightarrow{\mathrm{#1}}}
\providecommand{\vvt}[1]{\overrightarrow{\mathrm{#1}}}
\providecommand{\np}[2][]{{#2~#1}}
\providecommand{\vv}[1]{{\overrightarrow{#1}}}
\renewcommand{\d}{{\mathrm{d}}}
\providecommand{\degres}{{^\circ}}
\providecommand{\degre}{{^\circ}}
\providecommand{\blue}{\color{blue}}
\providecommand{\red}{\color{red}}
\providecommand{\enskip}{\quad}
\providecommand{\llbracket}{\lbrack\!\lbrack}
\providecommand{\rrbracket}{\rbrack\!\rbrack}
\providecommand{\euro}{{\text{€}}}
\providecommand{\no}{{\text{n°}}}
\providecommand{\interval}[2]{{[#1;#2]}}
\providecommand{\intervFO}[2]{{[#1;#2[}}
\providecommand{\dotfill}{\ldots}

\providecommand{\ds}{\displaystyle}
\providecommand{\textding}[1]{#1}
\providecommand{\Oij}{\(\mathrm{O};\vec{\imath},\vec{\jmath}\)}
\providecommand{\Oijk}{\(\mathrm{O};\vec{\imath},\vec{\jmath},\vec{k}\)}
\providecommand{\Ouv}{\(\mathrm{O};\vec{u},\vec{v}\)}
\providecommand{\intervFF}[2]{\ensuremath{[#1;#2]}}
\renewcommand{\euro}{\ifmmode\text{\texteuro}\else\texteuro\fi}
\DeclareUnicodeCharacter{00B0}{\ifmmode^{\circ}\else\textdegree\fi}
\DeclareUnicodeCharacter{20AC}{\ifmmode\text{\texteuro}\else\texteuro\fi}
% Énoncés : Éducation nationale. Transcription et corrigés : APMEP, Denis Vergès et contributeurs. Découpage et classement : CoopMaths, Sébastien Lozano, Éric Elter et contributeurs. Adaptation de la présentation : Maths Explorer.
% CC BY-SA : https://coopmaths.fr/www/about/
% https://coopmaths.fr/alea/static/eam/2026/tex/eam_specialite_2026_06_antilles_3_cor.tex
\begin{document}
\noindent\textit{APMEP / CoopMaths. CC BY-SA. Présentation adaptée par Maths Explorer.}\par\medskip
\begin{enumerate}
	\item \textbf{Affirmation 1 : Fausse.}

	Calculons les coordonnées des vecteurs $\vectt{AB}$ et $\vectt{AC}$.

	Les coordonnées de $\vectt{AB}$ sont \quad $\begin{pmatrix}x_{\mathrm{B}} - x_{\mathrm{A}}\\ y_{\mathrm{B}} - y_{\mathrm{A}} \end{pmatrix} $
	\quad soit, avec les valeurs données : \quad $\begin{pmatrix} 3 - 2\\2-0 \end{pmatrix} $.

	Finalement :\quad $\vectt{AB} \begin{pmatrix} 1\\2 \end{pmatrix} $.

	De même, on a :\quad $\vectt{AC} \begin{pmatrix} 0 - 2\\4 - 0 \end{pmatrix} $,\quad soit \quad $\vectt{AC} \begin{pmatrix} -2\\4 \end{pmatrix} $.

	Comme le repère est orthonormé, on peut utiliser la formule analytique du produit scalaire :

	$\vectt{AB}\cdot\vectt{AC} = x_{\vectt{AB}}x_{\vectt{AC}} + y_{\vectt{AB}}y_{\vectt{AC}} = 1 \times (-2) + 2 \times 4 = -2 + 8 = 6.$

	Le produit scalaire n'est pas nul, donc les vecteurs $\vectt{AB}$ et $\vectt{AC}$ ne sont pas orthogonaux, et donc les droites (AB) et (AC), que ces vecteurs dirigent, ne sont pas perpendiculaires.

	\item Déterminons l'expression de la fonction dérivée $f'$. $f$ est un produit, et donc on applique la formule pour dériver un produit.

On peut poser :\quad $u : x \longmapsto 3x - 1$ \quad et $v$ la fonction exponentielle.

On a alors :\quad $u': x \longmapsto 3$\quad et $v'$ également la fonction exponentielle (celle-ci étant, par définition, égale à sa dérivée).

On a alors :\quad $f = uv$\quad et donc \quad $f' = u'v + uv'$.

On en déduit, pour tout $x$ réel :\quad $f'(x) = 3\e^{x} + (3x - 1)\e^{x} = (3 + 3x - 1)\e^{x} = (3x + 2)\e^{x}$.

\begin{enumerate}
\item \textbf{Affirmation 2. a. : Fausse.}

En effet, la fonction exponentielle est à valeurs strictement positives, donc $f'(x)$ est du signe de $(3x + 2)$, mais cette expression change de signe.

$\begin{aligned}
3x + 2 > 0 &\iff 3x > -2\\
			&\iff x > -\dfrac{2}{3} \quad \text{ car on divise par }3>0
\end{aligned}$

On a donc le tableau de signes pour $f'$:

\hfill~\begin{tikzpicture}
\tkzTabInit[lgt=3]{$x$/1, signe de $f'(x)$/1}{$-\infty$,$-\dfrac{2}{3}$,$+\infty$}
\tkzTabLine{,-,z,+,}
\end{tikzpicture}\hfill~

On peut en déduire que la fonction $f$ est strictement décroissante sur l'intervalle $\left]-\infty ; -\dfrac{2}{3}\right] $, ce qui contredit que la fonction serait croissante sur cet intervalle. Ainsi la fonction n'est pas croissante sur $\R$ tout entier.

\item \textbf{Affirmation 2. b. : Vraie.}

On applique la formule pour obtenir une équation de la tangente :\quad $y = f'(a) (x - a) + f(a)$.

Ici, la fonction $f$ est celle de l'énoncé, et l'abscisse $a$ est égale à 0.

On va donc calculer :\quad $f(0) = (3\times 0 - 1)\e^{0} = -1 \times 1 = -1$\quad et \quad $f'(0) = (3\times 0 + 2)\e^{0} = 2\times 1 = 2$.

Ainsi, une équation de la tangente cherchée est :\quad $y = 2 \times (x - 0) + (-1)$.

Une fois simplifiée, cette équation réduite est bien :\quad $y = 2x - 1$.
	\end{enumerate}
\end{enumerate}

\end{document}
